% 第 1 章 About DC Explorer
\bichapter{关于 DC Explorer}{About DC Explorer}
\label{chap:about}

开发新 RTL，并将其与第三方 IP 及大量遗留 RTL 模块集成时，若缺少快速高效的手段来探索与改进数据、修复设计问题、并为 RTL 综合建立更好起点，往往非常耗时。
DC Explorer 通过支持早期 RTL 探索来解决这些问题，从而为 RTL 综合提供更好起点，并加速设计实现。

凭借对不完整设计数据的容错能力，以及显著快于完整综合的运行时间，DC Explorer 可尽早洞察实现结果；这些结果通常与 Design Compiler topographical 模式产出结果相差约在 10\% 以内。
工具使您能在设计周期早期高效进行多种设计配置的 what-if 分析，以加快高质量 RTL 与约束的开发，并推动更快、更易收敛的设计流程。
它还可生成早期网表，用于在 IC Compiler 中开始物理探索。通过一键访问 IC Compiler 设计规划，DC Explorer 允许您在设计周期极早期创建并修改 floorplan。

本章包括以下主题：
\begin{itemize}
  \item DC Explorer 设计流程
  \item 运行 DC Explorer 设计流程
  \item 使用 DC Explorer 的收益
  \item DC Explorer 与 Design Compiler 的兼容性
  \item IC Compiler 中的 Floorplan 探索
\end{itemize}

% ============================================================
\section{DC Explorer 设计流程}
\subsection*{DC Explorer Design Flow}

DC Explorer 设计流程提供快速综合，并帮助检测不完整与不匹配的数据。
DC Explorer 接受以下输入文件：RTL 网表、SDC 文件、逻辑库、物理库（可选）以及 floorplan（可选）。
工具生成报告、消息，以及可作为 IC Compiler 早期 floorplanning 输入的网表。

图~\ref{fig:dce-flow} 给出 DC Explorer 中的典型设计流程。

\figplaceholder{Figure 1: Design Flow in DC Explorer}{DC Explorer 中的设计流程}{fig:dce-flow}

\begin{seeAlsoBox}
运行 DC Explorer 设计流程，见下一节。
\end{seeAlsoBox}

% ============================================================
\section{运行 DC Explorer 设计流程}
\subsection*{Running the DC Explorer Design Flow}

按下列步骤运行 DC Explorer 设计流程：

\begin{enumerate}
  \item \textbf{创建 RTL 与约束。}\\
        DC Explorer 的输入设计文件为 RTL 描述、Synopsys Design Constraints（SDC）文件、逻辑库、物理库（可选）以及 floorplan（可选）。

  \item \textbf{打开设计。}\\
        可用以下任一方法打开设计：
        \begin{itemize}
          \item \cmd{analyze} 与 \cmd{elaborate} 命令
          \item \cmd{read} 类命令，例如 \cmd{read\_file}
        \end{itemize}

  \item \textbf{报告不完整或不匹配数据}，使用 \cmd{current\_design}、\cmd{link}、
        \cmd{check\_design} 与 \cmd{report\_design\_mismatch} 命令。\\
        在继续之前，您不一定需要先修复这些不完整或不匹配数据。

  \item \textbf{定义设计环境并施加设计约束。}
        \begin{itemize}
          \item 使用下列命令设置设计环境：\cmd{set\_operating\_conditions}、
                \cmd{set\_driving\_cell}、\cmd{set\_load} 与 \cmd{set\_fanout\_load}。
          \item 施加设计约束时使用 \cmd{read\_sdc} 命令。
        \end{itemize}

  \item \textbf{运行快速 compile 并优化设计}，使用 \cmd{compile\_exploration} 命令。

  \item \textbf{创建 HTML 格式的分类时序报告}，使用 \cmd{create\_qor\_snapshot} 与
        \cmd{query\_qor\_snapshot} 命令。\\
        检查报告中的缺失约束与不可行路径（infeasible paths）；继续之前可能需要修复这些不匹配。

  \item \textbf{写出网表}，对 \cmd{write\_file} 使用 \opt{-format} 选项。
\end{enumerate}

\begin{seeAlsoBox}
\begin{itemize}
  \item DC Explorer 设计流程（上文）
  \item 第~\ref{chap:tolerance}~章「对不完整或不匹配数据的容错」
  \item 第~\ref{chap:libs}~章「设置库」
  \item 第~\ref{chap:memory}~章「打开设计」
  \item 第~\ref{chap:env}~章「定义设计环境」
  \item 第~\ref{chap:constraints}~章「定义设计约束」
  \item 第~\ref{chap:opt}~章「优化」
  \item 第~\ref{chap:analyze}~章「报告结果质量」
  \item 第~\ref{chap:memory}~章「保存设计」
  \item 附录~\ref{app:filemgmt}「面向综合的设计文件管理」
\end{itemize}
\end{seeAlsoBox}

% ============================================================
\section{使用 DC Explorer 的收益}
\subsection*{Benefits of Using DC Explorer}

在 RTL 与约束开发阶段，DC Explorer 可通过早期快速综合提升效率。使用 DC Explorer 具有以下收益：
\begin{itemize}
  \item 产出快速综合结果
  \item 容忍不完整或不匹配的设计数据
  \item 在相同约束集下提供准确的初步综合，时序、面积与功耗数值通常与最终实现相差约在 10\% 以内
  \item 自动调和设计不匹配，并以简洁格式报告
  \item 在优化其他关键路径的同时检测不可行时序路径
  \item 生成便于导航的 HTML 时序报告，加速调试
  \item 创建可用于 IC Compiler floorplanning 的网表
\end{itemize}

% ============================================================
\section{DC Explorer 与 Design Compiler 的兼容性}
\subsection*{DC Explorer and Design Compiler Compatibility}

DC Explorer 在下列方面与 Design Compiler 兼容。默认情况下，DC Explorer：
\begin{itemize}
  \item 支持下列类别中的大多数 Design Compiler 命令：约束、脚本、报告、设计处理、设计输入/输出、建模与库设置
  \item 可运行既有 Design Compiler 脚本
  \item 对不支持的 Design Compiler 命令跳过执行并发出警告消息
\end{itemize}

要查看不支持的 Design Compiler 命令列表，请在 DC Explorer shell 提示符下输入 \cmd{help} 命令。

DC Explorer 与 Design Compiler 的差异如下：
\begin{itemize}
  \item DC Explorer 命令行界面（\cmd{de\_shell}）默认运行于 topographical 模式，因此无需指定 \opt{-topographical\_mode} 选项。
  \item 若尝试运行 \cmd{compile\_ultra}，DC Explorer 会改为调用 \cmd{compile\_exploration}。\\
        若在 \cmd{compile\_ultra} 或 \cmd{compile} 中使用不支持的选项，DC Explorer 会发出警告并跳过该选项。
        若使用 \opt{-incremental} 选项，则发出错误消息并完全跳过该命令的执行。
  \begin{itemize}
    \item DC Explorer 会忽略或转换不支持的 Design Compiler 命令，并继续执行脚本。
    \begin{itemize}
      \item 对设计或 QoR 影响较小的不支持命令会被忽略，并给出 DESH-009 警告。例如：
\begin{lstlisting}
Warning: Command 'set_fix_hold' is not supported in DC Explorer.
The command is ignored.(DESH-009)
\end{lstlisting}
      \item 影响显著的不支持命令会被忽略，并给出 DESH-008 错误。例如：
\begin{lstlisting}
Error: Command 'insert_dft' is not supported in DC Explorer.
(DESH-008)
\end{lstlisting}
      \item 自动化芯片综合（automated chip synthesis）命令不会被识别，DC Explorer 发出 CMD-005 错误消息。
    \end{itemize}
    \item DC Explorer 支持针对 pre-DFT 规则违例的测试设计规则检查（test design rule checking）。注意下列限制：
    \begin{itemize}
      \item 多数 DFT 规范命令可在已映射（mapped）设计上运行，但不能在未映射（unmapped）设计上运行。
      \item 不支持的命令会按影响程度以 DESH-009 警告或 DESH-008 错误忽略。
      \item 工具不支持 post-DFT DRC，即在已 stitch 的设计上运行 \cmd{dft\_drc} 命令。
    \end{itemize}
    有关 DFT 流程的更多信息，请参阅 TestMAX DFT 文档。
  \end{itemize}
\end{itemize}

\subsection{与 Design Compiler 的脚本兼容性}
\subsubsection*{Script Compatibility With Design Compiler}

在 DC Explorer 中运行 Design Compiler 脚本时，应遵循下列指南：
\begin{itemize}
  \item 若脚本检查程序名为 \cmd{dc\_shell} 及其提示符，需要将提示符从 \cmd{de\_shell} 改为 \cmd{dc\_shell}。
        为此，将变量 \varname{de\_rename\_shell\_name\_to\_dc\_shell} 设为 \texttt{true}。
        这会将只读环境变量 \varname{synopsys\_program\_name} 设为 \texttt{dc\_shell}。
  \item 若要编写仅在 DC Explorer（而非 Design Compiler）中执行某些命令的脚本，可使用如下结构：
\begin{lstlisting}
if {[shell_is_in_exploration_mode]} {do_DC_Explorer_setup}
\end{lstlisting}
  \item 要优化漏电功耗，请用 \cmd{set\_multi\_vth\_constraint} 替换 \cmd{set\_max\_leakage\_power}。
  \item 用一次 \cmd{compile\_exploration} 替换多次 \cmd{compile\_ultra}。
\end{itemize}

% ============================================================
\section{IC Compiler 中的 Floorplan 探索}
\subsection*{Floorplan Exploration in IC Compiler}

默认情况下，DC Explorer 使用 topographical 技术，使您能在 RTL 综合阶段准确预测布线后时序、面积与功耗，而无需基于线负载模型（wire load model）的时序近似。
它利用 Synopsys 的布局与优化技术，在综合内部驱动准确的时序预测，从而与最终物理设计具有更好相关性，并为物理实现提供更好的起点。

您可使用 DC Explorer 为 IC Compiler 的 floorplan 探索流程创建数据库。支持的数据库为 \texttt{.ddc} 或 ASCII Verilog 格式的已映射网表。
网表中多数不完整或不匹配数据会被解决，但逻辑库与物理库之间的不一致除外——在进入 IC Compiler floorplan 探索之前必须解决这些不一致。

从 DC Explorer 到 IC Compiler 设计规划的流程概览见原书插图；有关 GUI 信息，请参阅 \textit{Design Vision User Guide}。

\begin{seeAlsoBox}
\begin{itemize}
  \item 第~\ref{chap:gui}~章「执行 Floorplan 探索」
  \item 第~\ref{chap:floorplan}~章「使用 Floorplan 物理约束」
\end{itemize}
\end{seeAlsoBox}
